% RLJ main.tex Version 2025.0

\documentclass[10pt]{article} % For LaTeX2e

%%%%%%%%%%%%%%%%%%%%%%%%%%%%%%%%%%%%%%%%%%%%%%%%%%%%%%%%%%%%%%%%
% AUTHOR: Select ONE option:
%      [accepted]{rlj} --> for camera ready (after peer review, if accepted)
%      {rlj}           --> for submission
%      [preprint]{rlj} --> to de-anonymize and remove references to RLJ/RLC
%%%%%%%%%%%%%%%%%%%%%%%%%%%%%%%%%%%%%%%%%%%%%%%%%%%%%%%%%%%%%%%%
\usepackage{rlj}           % Should be uncommented for submission
%\usepackage[accepted]{rlj} % Should be uncommented for the camera-ready
%\usepackage[preprint]{rlj} % Should be uncommented for preprint versions

%%%%%%%%%%%%%%%%%%%%%%%%%%%%%%%%%%%%%%%%%%%%%%%%%%%%%%%%%%%%%%%%
% Recommended (but not required) packages
%%%%%%%%%%%%%%%%%%%%%%%%%%%%%%%%%%%%%%%%%%%%%%%%%%%%%%%%%%%%%%%%
\usepackage{mathtools}
\usepackage{mathrsfs}
\usepackage{graphicx}
\usepackage{subcaption}
\usepackage[space]{grffile}
\usepackage{url}
\usepackage{lipsum}
\usepackage{algorithm}
\usepackage{algpseudocode}
\usepackage{latexsym}
\usepackage{amssymb}
\usepackage{amsthm}
\usepackage{booktabs}
\usepackage{color}
\usepackage{enumitem}
\usepackage{float}
\usepackage{makecell}
\usepackage{microtype}
\usepackage{multirow}
\usepackage{xr}

%%%%%%%%%%%%%%%%%%%%%%%%%%%%%%%%%%%%%%%%%%%%%%%%%%%%%%%%%%%%%%%%
% AUTHOR: Fill in the following meta-data
%%%%%%%%%%%%%%%%%%%%%%%%%%%%%%%%%%%%%%%%%%%%%%%%%%%%%%%%%%%%%%%%

% Enter the title of your paper:
\title{RASPBERry: Enabling Continuous DRL at the Edge}
\setrunningtitle{RASPBERry: Enabling Continuous DRL at the Edge}

% WARNING: Authors must not appear in the submitted version. They should be hidden
% as long as the rlj package is used without the [accepted] or [preprint] options.
% Non-anonymous submissions will be rejected without review.

% Enter the author names below. 
% NOTE: Denote affiliations using superscripts as in the provided example.
% NOTE: Use \textscript{1,2,3} instead of $^{1,2,3}$.
%       - Failure to do so will cause affiliation superscripts to appear on the cover page for camera-ready and preprint versions.
\author{Chengming Zhou\textsuperscript{1,$\dagger$}, Stephen Mcgough\textsuperscript{1}, Jacek Cala\textsuperscript{1}}

% Author emails, which can be clustered if they have shared endings as in this example
\emails{\{c.zhou10, stephen.mcgough, jacek.cala\}@newcastle.ac.uk}

% Author affiliations, in the order the occur
% The inclusion of state/province, etc. is optional.
% The inclusion of multiple affiliations is optional.
%   - List multiple affiliations with comma-separated numbers as in the example.
\affiliations{$^{1}$\textbf{School of Computing Science, Newcastle University}}

% Define the summary that appears on the cover page.
\summary{
Deep Reinforcement Learning (DRL) often faces bottlenecks from long training cycles and high memory overheads for storing transitions.
We propose \textbf{RASPBERry} (RAM-Saving Prioritized Block Experience Replay), a novel replay-buffer architecture that compresses short sequences of consecutive transitions into blocks and assigns them a single priority. 
This design reduces sampling and update operations, slashing both memory demands and computational costs while preserving strong performance. 

We compared RASPBERry with Prioritized Experience Replay (PER) and Distributed PER (DPER) in a variety of environments, including Atari, Box2D, and MiniGrid.
Experiments on standard PCs, edge devices (comparable to an NVIDIA Jetson TX2), and HPC systems show that RASPBERry accelerates training substantially, often reaching final performance in roughly half the time needed by existing methods, and it reduce up to $97\%$ memory cost in experience replay.
It thus enables the continual off-policy DRL deployments on hardware with limited memory or compute capacity.
}


\contribution{
We provide a detailed analysis of the computational overhead in conventional PER and review how approaches such as Rainbow and Ray RLlib mitigate the costs of storing, sampling, and updating transitions under large mini-batch sizes.
}{None}

\contribution{
We propose \textbf{RASPBERry}, which compresses consecutive transitions into blocks and assigns a single priority per block, reducing replay operations and achieving up to a $97\%$ reduction in memory usage.
}{None}

\contribution{
We demonstrate RASPBERry’s effectiveness in experiments on local desktops, edge devices, and HPC clusters, confirming that its block-based design preserves scalability while lowering overhead in resource-constrained settings.
}{None}

% Include a list of keywords for the topic of the paper:
\keywords{Deep Reinforcement Learning, experience replay, memory optimization}



%%%%%%%%%%%%%%%%%%%%%%%%%%%%%%%%%%%%%%%%%%%%%%%%%%%%%%%%%%%%%%%%
%% Begin document, create title and abstract
%%%%%%%%%%%%%%%%%%%%%%%%%%%%%%%%%%%%%%%%%%%%%%%%%%%%%%%%%%%%%%%%
\begin{document}

\makeCover  % Create the cover page
\maketitle  % Make the title section

\begin{abstract}
In the field of deep reinforcement learning, long training times and high DRAM costs have been significant challenges.
We introduce a novel method, RASPBERry (RAM-Saving Prioritized Block Experience Replay), along with its distributed version, which reorganizes the structure of the replay buffer to drastically enhance the speed of buffer operations while significantly reducing DRAM requirements.
By grouping consecutive transitions into compressed blocks and assigning a unified priority to each block, RASPBERry reduces the number of sampling and update operations, thereby accelerating the overall training process.

We demonstrate the effectiveness of RASPBERry by integrating it with Prioritized Experience Replay (PER) and Distributed PER (DPER) across a variety of environments, including Atari, Box2D, and MiniGrid.
Our approach significantly improves training speed and efficiency: for some Atari games, RASPBERry achieves performance comparable to DPER in roughly $50\%$ of the training time when using a standard PC—far fewer computational resources than required by the original implementation.
Specially, on an 8-core laptop with poor GPU, we solve Pong in one hour and Boxing in two hours.
Moreover, we simulate an edge device environment—comparable to an NVIDIA Jetson TX2 equipped with only 8 GB of DRAM and modest CPU/GPU power—to show that DDQN with RASPBERry can solve Pong within 8 hours on edge.

This achievement demonstrates that RASPBERry enables Deep Q-learning algorithms with experience replay to operate efficiently on resource-constrained systems include richer state representation, paving the way for their application in edge devices such as robotic vacuum cleaners and autonomous vehicles.

\end{abstract}

\section{Introduction}\label{sec:intro}
Deep Reinforcement Learning (DRL) integrates reinforcement learning (RL) with deep neural networks, achieving strong performance in Atari games \citep{HumanlevelControlDeep_mnih_2015}, robotics, and continuous control.
A foundational element in many DRL algorithms, such as deep Q-networks (DQN), is \emph{Experience Replay} (ER) \citep{SelfimprovingReactiveAgents_lin_1992}. 
Building on ER, Prioritized Experience Replay (PER) \citep{GeneralizedPrioritizedSweeping_andre_1997, PrioritizedExperienceReplay_schaul_2016} samples transitions by their Temporal Difference (TD) error, a proxy for each transition's surprise or importance.
However, large replay buffers with high-dimensional inputs can incur significant computational and storing overhead.  
In particular, when training with large buffers and high mini-batch sizes, the overhead during \texttt{Storing}, \texttt{Sampling}, and \texttt{Updating} transitions becomes a bottleneck, and memory usage can balloon to tens or hundreds of gigabytes---challenging for resource-constrained scenarios.  
Although optimizations like fixed-size tensor arrays \citep{CurtParkRainbowisallyouneed_parkcurt_2024, StableBaselines3ReliableReinforcement_raffin_2021} can help, they often lack flexibility for dynamic buffers or large-scale distributed setups \citep{RayDistributedFramework_moritz_2018}. 
Other optimizations, like asynchronous or distributed replay shards \citep{RayDistributedFramework_moritz_2018}, can exceed many practitioners' available resources.

We propose \textbf{RASPBERry} (RAM-Saving Prioritized Block Experience Replay), a novel method that reorganizes the replay buffer to cut both computational and memory overhead.
Instead of storing and sampling individual transitions, RASPBERry compresses short consecutive transition sequences into blocks, assigning them a single priority.
This grouping lowers the costs of sampling and updates while exploiting redundancy in consecutive frames for compression.
Experiments show that RASPBERry reduces memory usage by up to $97\%$ while maintaining strong performance on several benchmarks.
For instance, it stores one million Atari transitions using just 4\,GB of DRAM (versus about 57\,GB for standard PER).  
On a simulated edge device (similar to an NVIDIA Jetson TX2 \citep{JetsonTX2Module_nvidiadeveloper_2024} with 8\,GB RAM), RASPBERry allows Double DQN (DDQN) to solve Atari Pong and Boxing in less than 8 hours.
This advancement paves the way for deploying advanced off-policy DRL algorithms on edge devices such as sweeping robot, Drone grand control station and other low-resource platforms.
The code are available in \textbf{hidden GitHub}.

The contributions of this paper are as follows: 
\begin{enumerate} 
    \item We provide a detailed analysis of the computational overhead in conventional PER, identifying how large mini-batch sizes exacerbate the cost of storing, sampling, and updating transitions. 
    \item We present RASPBERry, which groups consecutive transitions into compressed blocks, drastically reducing buffer manipulation costs and memory requirements.
    \item We demonstrate that RASPBERry can be integrated with existing PER-based frameworks, achieving comparable or superior training performance while drastically lowering resource demands. 
\end{enumerate}

%%%%%%%%%%%%%%%%%%%%%%%%%%%%%%%%%%%%%%%%%%%%%%%%%%%%%%%%%%%%%%%%
%% Section: Analysis
%%%%%%%%%%%%%%%%%%%%%%%%%%%%%%%%%%%%%%%%%%%%%%%%%%%%%%%%%%%%%%%%

\textcolor{red}{Discussion on Experience Replay Algorithm}

\begin{itemize}
    \item Batch Size for Training RL. 
    \item Training intensity. 
    \item Size of experience replay. 
    \item Single-node or multi-node. DQN, Ape-X DDQN. 
    \item Operations in Prioritized Experience Replay.  (Sample, Update). 
\end{itemize}

\textcolor{red}{Discussion on time and RAM for current RL algorithm. }

\section{Analysis of computational load in PER}
\label{sec:motivation}



We first measure how sampling and vectorization affect PER in DDQN \citep{PrioritizedExperienceReplay_schaul_2016} under varying batch sizes.  
We run four Atari games (Beam-Rider, Breakout, Qbert, Space-Invaders) on Ray \citep{RayDistributedFramework_moritm shagnz_2018} using a single replay buffer thread.

\begin{figure}[t]
    \centering
    \begin{minipage}[b]{0.48\linewidth}
        \centering
        \includegraphics[height=2.4cm]{images/motivation/stage_stacks_0115.png}
        \caption{DQN Training Time Breakdown}
        \label{fig:operation_stack}
    \end{minipage}
    \hfill
    \begin{minipage}[b]{0.48\linewidth}
        \centering
        \includegraphics[height=2.4cm]{images/motivation/operation_line_0115.png}
        \caption{Buffer operations broken down}
        \label{fig:buffer_operation_breakdown}
    \end{minipage}
\end{figure}

\textcolor{red}{Add discussion on ideal situation, and we can add our method's improvement in this part. }

Figure~\ref{fig:operation_stack} shows how total training time is split among acquiring new experiences (\textit{Observe}), neural network training (\textit{Train}), and replay buffer operations (\textit{Buffer}).  
As the mini-batch size increases, buffer manipulation becomes a bigger fraction of the total time.
Figure~\ref{fig:buffer_operation_breakdown} further breaks down buffer overhead into \texttt{Store} ($T_{st}$), \texttt{Sample} ($T_{sa}$), and \texttt{Update} ($T_{u}$). 
Specifically, \texttt{Store} measures how long it takes to add new transitions with default TD errors, \texttt{Sample} covers selecting transitions from the replay buffer (including tensor conversion and GPU transfer), and \texttt{Update} recalculates and inserts updated TD errors into the sum-tree structure.  
Separate time measurements for each buffer operation, depicted in Figure~\ref{fig:operation_stack}, reveal a clear trend, \texttt{Sample} and \texttt{Update} grow significantly with larger batch sizes while \texttt{Store} remains relatively stable.

Our contribution:
\textcolor{red}{We are the first to dramatically reduce the RAM consumption of PER algorithm,}

\begin{figure}[t!h]
  \centering
  % Left figure block
  \begin{minipage}[c]{0.48\linewidth}
    \centering
    \includegraphics[height=2.4cm]{images/motivation/mean_n_processed_transitions_per_0115.png}
    \captionof{figure}{Processed frames per in DQN}
    \label{fig:per_n_processed}
  \end{minipage}
  \hfill
  % Right figure block
  \begin{minipage}[c]{0.48\linewidth}
    \centering
    \includegraphics[height=2.4cm]{images/motivation/mean_n_processed_transitions_dper_0115.png}
    \captionof{figure}{Processed frames per in Ape-X DDQN}
    \label{fig:dper_n_processed}
  \end{minipage}
\end{figure}

The inefficiencies in these buffer operations can become a bottleneck in the DDQN training cycle.
Figures~\ref{fig:per_n_processed} and~\ref{fig:dper_n_processed} show transitions processed per second under DQN+PER and Ape-X DDQN+DPER \citep{DistributedPrioritizedExperience_horgan_2018} under different batch. 
The training intensity of Ape-X DDQN+DPER scaled in proportion to the batch size. 
When the mini-batch size exceeds 256 in \textbf{PER} or 512 in \textbf{DPER}, throughput stops scaling linearly and even decreases.  
Hence, buffer overhead forms a bottleneck across both single and distributed settings, limiting the benefits of high batch sizes.

\section{Related Work}
Neuroscience suggests that replay in the mammalian hippocampus often unfolds as consecutive memory sequences \citep{ReverseReplayBehavioural_foster_2006}, potentially boosting memory consolidation or planning \citep{HippocampalPlaceCells_olafsdottir_2015}, and supporting memory functions like preparing for future tasks \citep{TheoryHippocampalFunction_rolls_2024}.
In DRL, \citet{HumanlevelControlDeep_mnih_2015} introduced ER to reduce correlated updates in DQN, and \citet{PrioritizedExperienceReplay_schaul_2016} added TD-error-based prioritization to focus on more informative transitions.
Many works improve sample efficiency via other sampling or selection heuristics: for example, \citet{PrioritizedSequenceExperience_brittain_2020} up-weights recent transitions leading to large TD errors,
\citet{SampleEfficientDeepReinforcement_lee_2019} propose backward updates from entire episodes, and \citet{SelfOrganizingMapsStorage_karimpanal_2018} select transition subsequences for faster learning.
However, these methods still operate at the individual transition level, leaving replay overhead largely unaddressed.

Meanwhile, distributed methods like Ape-X \citep{DistributedPrioritizedExperience_horgan_2018} reduce \texttt{Observe} overhead with multiple actors but can face communication and tuning complexities.
Frameworks such as Ray RLlib \citep{RayDistributedFramework_moritz_2018} boost throughput with multiple replay buffers.
Rainbow implementation solve this issue by fixed-size tensor arrays \citep{CurtParkRainbowisallyouneed_parkcurt_2024}, but this implementation lack adaptability for dynamic scaling \citep{RLx2TrainingSparse_tan_2023} and other replay researches.

Drawing on insights that replay often occurs in sequences \citep{ReverseReplayBehavioural_foster_2006, ReplayCorticalSpiking_vaz_2020} and that PER focuses on sampling salient transitions \citep{PrioritizedExperienceReplay_schaul_2016}, our approach instead uses \emph{short transition blocks} as the fundamental replay unit.  
Unlike Rainbow’s rigid tensor approach \citep{CurtParkRainbowisallyouneed_parkcurt_2024}, our method accommodates dynamic buffer resizing and achieves up to 97\% reduction in DRAM usage.
This design also integrates seamlessly with distributed setups, such as Ape-X’s multi-actor architecture \citep{DistributedPrioritizedExperience_horgan_2018} and Ray’s multi-buffer framework \citep{RayprojectRlexperiments_ray_2022}. 
Our experiments confirm, converting transitions into short, compressed blocks not only avoids high replay overhead but also preserves training performance, indicating that block-level sampling and updates can effectively merge the benefits of sequential replay with practical memory savings.

\section{RAM Saving Prioritised Block Experience Replay (RASPBERry)}\label{sec:contribution}
We propose RASPBERry to address the twin challenges of high replay overhead and heavy DRAM requirements in ER by compressing and storing \textbf{blocks} of transitions instead of \textbf{individual} transitions.

RASPBERry adopts a dual-buffer structure: 
\begin{itemize}[leftmargin=*]
    \item $\mathcal{D}_{\text{block}}$: a small buffer accumulating new transitions.
    Once it has enough data, these transitions are grouped into a single block and transferred to $\mathcal{D}$.
    \item $\mathcal{D}$: the main replay buffer, storing each block as one element: $(B_{S_{t-1}},\, B_{A_{t-1}},\, B_{R_t},\, B_{\gamma_t},\, B_{S_t})$, each of size $m$ transitions, akin to the stacking strategy of Rainbow \citep{RainbowCombiningImprovements_hessel_2017}.
\end{itemize}

By treating each block in $\mathcal{D}$ as a single element during buffer operations, RASPBERry dramatically reduces the required sampling and update operations.
For instance, if mini-batch size $M=512$ and block size $m=32$, only $512/32=16$ sampling operations are performed rather than 512.
Furthermore, blocks also compress well when observations are consecutive, cutting memory usage.
For Atari, a $2\%$ final size is possible under typical Deflate-based~\citep{DEFLATECompressedData_deutsch_1996} compression.
Such savings enable, for example, storing one million transitions in about 4\,GB of DRAM, compared to 56.5\,GB under other replay buffer architectures.

Figure~\ref{fig:Diff} compares a single training iteration in PER (a) and RASPBERry (b), where the mini-batch contains $M = 4$ transitions, $\mathcal{D}_{\text{block}}$ holds $m = 2$ transitions.
Instead of one TD-error per transition, RASPBERry stores a single priority for the entire block, updated by the \textbf{average} TD-error of those transitions.  
This design lowers sum-tree update overhead while retaining the benefits of prioritized replay; meanwhile, buffer $\mathcal{D}$ keeps a flexible structure (e.g. an ordered list with a sum tree), allowing easy integration with dynamic buffer strategies.
It should be noted that we store a small number (tens of blocks) of data for parallel compression, which results in storing new experiences in \textbf{pulses}, which may affect sampling efficiency.

\begin{figure}[ht]
    \centering
    \includegraphics[width=0.9\textwidth]{images/workflow-PER-RS.jpg}
    \caption{Single training iteration in PER (a) vs. RASPBERry (b).}
    \label{fig:Diff}
\end{figure}

To enable RASPBERry in distributed frameworks such as Ape-X \citep{DistributedPrioritizedExperience_horgan_2018}, we modify the Replay and Learner components to support block-level sampling and updates, leaving the Actor component unchanged.  
Pseudocode for the Replay and Learner processes is provided in Appendix~\ref{apx:raspberry_code}, highlighting the changes from the standard Ape-X architecture.

Empirically, RASPBERry lowers buffer operation time and memory use.
For example, with a replay capacity of 500k transitions, a single-Actor Ape-X DDQN can run with just 8\,GB of memory instead of the 70\,GB required by DPER---vital for resource-constrained platforms.

\section{Experiments}
\label{sec:experiments}
We evaluated RASPBERry vs.\ standard PER and DPER using Ray-Project RLLib v2.8 \citep{RayprojectRlexperiments_ray_2022}, reproducing benchmark configurations.
Our experiments spanned multiple Atari games, as well as Box2D and MiniGrid \citep{MinigridMiniworldModular_chevalier-boisvert_2023, Gymnasium_towers_2024} environments, to assess the robustness and scalability of the proposed methods.
We used two hardware setups: (i) a local PC, and (ii) a high-performance computing (HPC) system. 
Detailed configurations of the hardware and the hyperparameters for DDQN and Ape-X experiments are provided in Appendix~\ref{apx:experiments}.
Although our results cannot be directly compared with those of~\citet{PrioritizedExperienceReplay_schaul_2016} due to differences in hyperparameters, environment versions, and hardware, we have controlled all other factors except for the replay buffer and reproduced experiments on our own hardware/software stack.
The relative performance gains thus remain valid under our experimental settings.

\begin{table*}[t]
\centering\small
\begin{tabular}{lccccccc}
\toprule
 & \multicolumn{3}{c}{PER vs RASPBERry} & \multicolumn{3}{c}{DPER vs RASPBERry} & Compress  \\
\midrule
Environment & $R_{\text{improve}}$ & $T_{\text{save}}$ & $R_{\text{max}}$ & $R_{\text{improve}}$ & $T_{\text{save}}$ & $R_{\text{max}}$ & $C_{\text{compress}}$\\
\midrule
BeamRider        & 56.3\%  & 72.2\%  & 66.7\%  & 14.4\%  & 88.1\%  & 146.8\%  &  89\% \\
Breakout        & 80.2\% & 82.1\% & 93.9\% & 201.6\% & 96.8\% & 104.0\% &  88\% \\
Qbert            & 67.5\%  & 45.4\%  & 124.0\% & 116.9\% & 77.9\%  & 102.3\%  &  89\% \\
Pong*             & 64.1\%  & 85.8\%  & 100.0\% & 117.3\% & 82.6\%  & 100.0\%  &  95\% \\
SpaceInvaders    & 8.4\%   & 27.5\%  & 110.3\% & 23.5\%  & 56.2\%  & 89.6\%   &  89\% \\
LavaCrossingS9N1* & 260.8\% & 82.1\%  & 100.0\% & 10.2\%  & 42.1\%  & 100.0\%  &  93\% \\
Empty-8x8*       & 15.2\%  & 0.0\%   & 100.0\% & 31.2\%  & 60.0\%  & 99.9\%   &  97\% \\
CarRacing*        & -36.9\% & -47.1\% & 99.3\%  & 10.5\%  & 83.3\%  & 100.0\%  &  75\% \\
\midrule
Trimmed Mean     & 48.6\%  & 51.6\%  & 100.6\% & 52.3\%  & 74.7\%  & 101\%  &  89.2\% \\
\bottomrule
\end{tabular}
\caption{Summary of Local Experiment Results, Including Compression Rates
\newline \textbf{*} Environments marked with `*' have a known maximum episodic reward}
\label{tab:local_game_summary}
\end{table*}

\subsection{RASPBERry vs PER} \label{sec:diff_per_pber}
We first compare DDQN+PER to DDQN+RASPBERry with block size $m=8$.
Hence, sampling $M'=4$ blocks covers $32$ transitions effectively (the default batch size in Ray ~\citep{RayprojectRlexperiments_ray_2022}).  

\begin{figure}[b]
    \centering
    \begin{minipage}[b]{0.34\linewidth}
      \centering
      \includegraphics[height=0.35\linewidth]{images/analysis/RASPBERry_time_usage.png}
      \caption{replay buffer operations time breakdown}
      \label{fig:local_per_time_breakdown}
    \end{minipage}
    \hfill
    \begin{minipage}[b]{0.28\linewidth}
      \centering
      \includegraphics[height=0.55\linewidth]{images/t_save.png}
      \caption{$T_{\text{save}}$ illustration, $t_2$ is the medium}
      \label{fig:t_save}
    \end{minipage}
    \hfill
    \begin{minipage}[b]{0.34\linewidth}
      \centering
      \includegraphics[height=0.38\linewidth]{images/analysis/edge_lr.png}
      \caption{Learning curves of Atari Pong on Edge}
      \label{fig:lr_edge_pong}
    \end{minipage}
\end{figure}

We focusing on those key metrics: $T_{\text{ER}}$, $R_{\text{improve}}$, $T_{\text{save}}$,  $R_{\text{max}}$ and $C_{\text{compress}}$.
First, we analysed the average time breakdown for PER and RASPBERry (Figure~\ref{fig:local_per_time_breakdown}), showing a $45\%$ average reduction in $T_{\text{ER}}$.  
This time savings allows DDQN with RASPBERry to process $22-84\%$ more frames in the same `wall-clock time', significantly improving training throughput.
Then, we calculated $R_{\text{improve}}$, $T_{\text{save}}$, $R_{\text{max}}$ and $C_{\text{compress}}$ for each environment (Table~\ref{tab:local_game_summary}), showing that RASPBERry consistently improves performance compared to PER.
\textbf{$R_{\text{improve}}$} is the median relative reward improvement in a 15-minute rolling window, capturing incremental gains (Eq.~\eqref{eq:er_reward}).
\textbf{$T_{\text{save}}$} is the median fraction of time saved in reaching $50$ predefined reward thresholds (Eq.~\eqref{eq:er_save}), measured either in wall-clock time or in frames. 
Figure~\ref{fig:t_save} illustrates how $T_{\text{save}}$ is computed via hypothetical learning curves partitioned into five score points ($R_1,\dots,R_5$).
For each $R_n$, $t_{\text{PER}}^{(n)}$ is the time for DDQN with PER to hit $R_n$, and $t_{\text{RASPBERry}}^{(n)}$ is the time for RASPBERry.  
\textbf{$R_{\text{max}}$} is each method’s highest episodic reward throughout training, reflecting ultimate achievable performance.
For instance, on Qbert, RASPBERry increases the reward by $67.5\%$ within the same `wall clock time', saves $72.2\%$ of the time needed to reach comparable reward, and achieves a $124\%$ maximum reward compared to the baseline. 
Finally, the compression rate $C_{\text{compress}}$ is also around 89\%, indicating significant memory savings, allowing for more efficient replay buffer management.
We also confirmed those performance gains by experiments on a HPC system; see section \ref{sec:extra_hardware}.

\begin{equation}
R_{\text{improve}} = \mathrm{median}(\frac{R_{\text{RASPBERry, window}} - R_{\text{PER, window}}}{R_{\text{PER, window}}})
\label{eq:er_reward}
\end{equation}
\begin{equation} \label{eq:er_save}
T_{\text{save}} = \mathrm{median}\left(\left\{\frac{t_{\text{PER}}^{(n)} - t_{\text{RASPBERry}}^{(n)}}{t_{\text{PER}}^{(n)}}\right\}_{n=1}^{N}\right)
\end{equation}

\begin{figure*}[t]
    \centering\includegraphics[width=0.9\textwidth]{images/experiments/all_local_per_time_10_paper.png}
    \caption{Learning curves of DDQN experiments, ran repeated}
    \label{fig:local_per}
\end{figure*}

\subsection{distributed RASPBERry vs DPER} \label{sec:diff_dper_dpber}
We repeat the above comparison under Ape-X DDQN+DPER.  
As in the previous experiments, the mini-batch size is $512$, and we use a block size of $m = 16$ for distributed RASPBERry, with $M' = 32$ blocks sampled per iteration.  
Distributed RASPBERry reduces observation speed by $42\%$ but processes $1100\%$ more frames, resulting in a $R_{\text{improve}}$ of $52.3\%$ and a $T_{\text{save}}$ of $74.7\%$.

\begin{figure*}[ht]
    \centering
    \includegraphics[width=0.9\textwidth]{images/experiments/all_local_dper_time_10_paper.png}
    \caption{Learning curves of Ape-X DDQN experiments, ran repeated.}
    \label{fig:local_dper}
\end{figure*}

\subsection{Experiments on edge and HPC} \label{sec:extra_hardware}
Our experiments on edge devices and a HPC system validate the practical deployment of RASPBERry in both resource-constrained and large-scale environments.

\subsubsection{Edge Deployment.}
One of the key advantages of RASPBERry is its ability to run efficiently on edge devices with limited memory and processing power.  
We conducted experiments using simulated edge hardware environments, similar to the NVIDIA Jetson TX2, which is constrained by 8\,GB of DRAM and a quarter of an NVIDIA P4 GPU.  
Figure~\ref{fig:lr_edge_pong} depicts the learning curves for Pong, showing that RASPBERry can effectively handle a replay buffer of up to one million transitions on these edge devices, whereas PER can only store up to 90k transitions under the same memory limit.  
RASPBERry’s memory efficiency enables better utilization of available resources, allowing it to achieve faster training and improved performance in comparison to the standard PER approach.

The ability to store more transitions with RASPBERry, while maintaining efficient processing, makes it particularly suitable for edge environments where memory and compute resources are often limited.  
This is critical for deploying of-policy DRL algorithms in real-world applications, such as robotics or autonomous vehicles, where edge devices are commonly used for on-device learning.

\subsubsection{HPC Validation.}
To further validate the scalability of RASPBERry, we conducted large-scale experiments on a high-performance computing (HPC) system.
We ran experiments across 24 Atari games to test RASPBERry’s performance in a multi-node, multi-agent environment, simulating large-scale deployments of DRL algorithms.  
While we could not replicate the exact environment and configuration of prior studies, we ensured the generalizability of our results by maintaining consistent hyperparameters and environmental setups.

Table~\ref{tab:hpc_sum} shows that RASPBERry outperforms the baseline (PER/DPER) in terms of both reward improvement and time savings across various games.  
In particular, RASPBERry processes $16--96\%$ more frames compared to PER, demonstrating its ability to handle larger workloads more efficiently in a distributed setting.  
We hypothesize that the compressed replay buffer used in RASPBERry facilitates better memory addressing and cache utilization, leading to faster sampling and training times, even on large DRAM systems.
The results from our HPC experiments confirm that RASPBERry’s performance advantages are not limited to edge devices but extend to large-scale systems, making it a versatile solution for different deployment scenarios.  

\begin{table}[h]
% \footnotesize
\centering
\begin{tabular}{lrrrrr}
\toprule
 & $R_{\text{improve}}$ & $T_{\text{save}}$ & $R_{\text{max}}$ & \textbf{Sampled} & \textbf{Trained} \\
\midrule
DDQN + RASPBERry       &  60.4\% & 35.4\% & 102.9\% & 160\% & 160\% \\
Ape-X DDQN + RASPBERry &  78.6\% & 56.5\% & 122.7\% &  35\% & 1035\% \\
\bottomrule
\end{tabular}
\caption{High-level comparison of RASPBERry vs.\ standard PER/DPER in HPC experiments. Baseline $R_{\text{max}}$ are 100\%.}
\label{tab:hpc_sum}
\end{table}

Overall, the experiments demonstrate that RASPBERry is not only well-suited for edge environments where memory and computational power are limited but also scalable in large HPC systems.  
This highlights its versatility and potential for deployment in a wide range of applications, from edge computing to large-scale distributed systems.
% What is more, it is foreseeable that this optimization will also be effective for other ER methods such as HER \citep{HindsightExperienceReplay_andrychowicz_2018}, or PER in R2D2 \citep{R2D2RepeatableReliable_revaud_2019}.

\section{Analysis}
We investigate the mechanisms behind RASPBERry's performance gains and its broader applicability.
\subsection{Where does the improvement come from?}
\label{sec:where_come_from}

Our results in Sections~\ref{sec:diff_per_pber} and \ref{sec:diff_dper_dpber} indicate that RASPBERry consistently enhances training performance compared to traditional PER/DPER.  
To pinpoint the source of these improvements, we analyse performance in terms of the number of \emph{trained} and \emph{sampled} frames.  
This perspective allows us to distinguish whether the gains stem primarily from faster processing, better sample efficiency, or a combination of both.

Table~\ref{tab:performance_comparison} summarizes local experiments, contrasting the metrics when normalized by the number of \emph{trained frames} vs.\ \emph{sampled frames}.
We observe that DDQN with RASPBERry exhibits better scores under both perspectives, indicating that the \emph{pulsed storing} of transitions can yield higher sample efficiency.  
For Ape-X DDQN, the \textbf{Trained Frame} row shows a value of $-170\%$ in $T_{\text{save}}$, which at first might seem counter-intuitive. 
In fact, RASPBERry processes more transitions in the same time span, effectively distorting this particular metric.
The overhead reduction in buffer operations repositions the bottleneck in the Ape-X pipeline, thereby yielding an overall net gain in throughput despite a local slowdown in \textbf{Actor–Replay} communication.

\begin{table}[ht]
\centering\small
\setlength{\tabcolsep}{4pt}
\begin{tabular}{lrrrrrr}
\toprule
\multirow{2}{*}{} & \multicolumn{3}{c}{Trained Frame} & \multicolumn{3}{c}{Sampled Frame} \\
& $T_{\text{save}}$ & $R_{\text{improve}}$ & $R_{\text{max}}$ & $T_{\text{save}}$ & $R_{\text{improve}}$ & $R_{\text{max}}$ \\
\midrule
DQN + RASPBERry & 30\% & 46\% & 71\% & 30\% & 45\% & 71\% \\
Ape-X DDQN + RASPBERry & -170\% & -40\% & 14\% & 89\% & 35\% & 14\% \\
\bottomrule
\end{tabular}
\caption{Comparison across local experiments. 
Negative ones indicate slower or lower performance under that specific viewpoint.}
\label{tab:performance_comparison}
\end{table}

Overall, these results suggest that RASPBERry’s main advantage lies in faster replay buffer operations, allowing up to 1100\% more frames to be processed in Ape-X.  
Figure~\ref{fig:selected_by_frames} illustrates learning curves from selected experiments, plotted by the number of trained frames. 
We note that RASPBERry often outperforms PER in early training stages, hinting that batching newly arrived transitions into blocks accelerates learning initially.

\begin{figure}[ht]
    \centering
    \includegraphics[width=0.9\linewidth]{images/analysis/selected_by_frames_time_paper.png}
    \caption{Learning curves of selected experiments by trained frames.}
    \label{fig:selected_by_frames}
\end{figure}

\subsection{How does block size affect the training?}
We further examined the role of the block size $\mathcal{D}_{\text{batch}}$ by varying $m$ in RASPBERry. 
Figure~\ref{fig:bz_analysis_dper} shows learning curves for selected Atari games under DPER, compared to different block sizes in RASPBERry.
The performance benefit declines notably once $m$ exceeds $\sqrt{M}$, where $M$ is the mini-batch size.  
This suggests that $m \approx \sqrt{M}$ might be an effective heuristic, though more exhaustive experimentation and theoretical analysis are needed to confirm if $\sqrt{M}$ is optimal in broader scenarios.

\begin{figure}[ht]
    \centering
    \includegraphics[width=0.8\linewidth]{images/analysis/analysis_dpber_diff_bz.png}
    \caption{Varying block size in RASPBERry on selected Atari games, compared with DPER.}
    \label{fig:bz_analysis_dper}
\end{figure}

Such preliminary findings highlight a trade-off between reducing replay overhead and maintaining sufficient data granularity.
In practice, a moderate block size appears to balance computational efficiency with data diversity, though we leave a more rigorous study of this balance to future work.


\section{Conclusion and Future Work}
We have presented \textbf{RAM Saving Prioritised Block Experience Replay} (RASPBERry), a novel approach that restructures and manages the replay buffer to alleviate computational bottlenecks and reduce memory demands in off-policy DRL.  
Across diverse hardware platforms---from local PCs, edge to HPC clusters---RASPBERry consistently demonstrated superior training throughput and memory efficiency compared to standard PER-based methods.
Notably, it achieves up to a 97\% reduction in DRAM usage, enabling off-policy algorithms to run on memory-constrained  devices.

Unlike the fixed-size tensor strategy employed in methods such as Rainbow \citep{CurtParkRainbowisallyouneed_parkcurt_2024}, RASPBERry maintains broader compatibility by storing transitions in compressed blocks while retaining a flexible storing structure.  
This flexibility not only eases deployment in large-scale environments, but also opens avenues for further improvements.  
For instance, if future versions of the Ray framework were to directly handle binary blocks in the replay buffer, even higher compression ratios might be achieved, effectively minimizing the DRAM footprint.  

In future work, exploring more sophisticated weighting and sampling strategies within each block would be a interesting topic.
Techniques like Batch Prioritized Experience Replay \citep{OffPolicyCorrectionDeep_cicek_2021} could be incorporated to rebalance the priorities of transitions at a finer granularity, thereby enhancing training stability and performance.  
We believe that combining RASPBERry’s block-based replay with these advanced weighting schemes will lead to even more efficient and scalable DRL solutions, further extending the applicability of off-policy methods in resource-constrained and large-scale settings.

\newpage

%%%%%%%%%%%%%%%%%%%%%%%%%%%%%%%%%%%%%%%%%%%%%%%%%%%%%%%%%%%%%%%%
%% Appendices
%%%%%%%%%%%%%%%%%%%%%%%%%%%%%%%%%%%%%%%%%%%%%%%%%%%%%%%%%%%%%%%%
\appendix

\section{RASPBERry Implementation in Ape-X architecture}
\label{apx:raspberry_code}
This section provides the pseudocode for integrating \textbf{RASPBERry} into the Ape-X DDQN architecture \citep{DistributedPrioritizedExperience_horgan_2018}.
While Section~\ref{sec:contribution} in the main text offers an overview, here we present the complete pseudocode for both the \emph{Replay} (Algorithm~\ref{alg:replay}) and the \emph{Learner} (Algorithm~\ref{alg:learner}) processes.  
Modifications or new procedures appear in \textcolor{blue}{blue} for clarity.

\subsection{Key Differences from Standard Ape-X.}
RASPBERry introduces:
\begin{itemize}[leftmargin=*]
    \item \emph{Block-level storage:} Transitions arrive from actors, are grouped into short blocks, and compressed before being stored in the replay buffer.
    \item \emph{Single-priority updates per block:} Sum-tree operations occur once per block rather than per transition, drastically cutting overhead.
    \item \emph{Flexible block decomposition:} During sampling, blocks are decompressed to reconstruct the original transitions needed for training.
\end{itemize}

In Algorithm~\ref{alg:replay}, newly arrived transitions are split into blocks (\textcolor{blue}{line~5}), compressed, and stored as block-level entities.
In Algorithm~\ref{alg:learner}, each sampled block is decompressed (\textcolor{blue}{line~4}) prior to network training; updated priorities are aggregated and communicated back (\textcolor{blue}{lines~8--9}). 
These modifications allow RASPBERry to reduce sampling and update calls from $M$ per mini-batch to $M/m$, where $m$ is the block size.  
For more details on block compression rates, sum-tree operations, and memory consumption, see Section~\ref{sec:contribution}.

\begin{algorithm}[ht]
\caption{RASPBERry Replay in Ape-X}
\label{alg:replay}
\small{
\begin{algorithmic}[1]
\Procedure{Replay}{$T$}\Comment{Store and sample transitions for $T$ timesteps}
    \State $R_0 \gets \Call{Initialize Replay Buffer}{}$ 
    \For{$t = 1$ {\bfseries to} $T$}
        \State $\tau, p \gets \Call{replay.get}{\,}$ 
        \Comment{Get a batch of transitions $\tau$ and corresponding priorities $p$ from Actors}
        \State \textcolor{blue}{$B_{\tau}, B_{p} \gets \Call{SplitIntoBlock}{\tau, p}$}
        \Comment{Group transitions into blocks, compress and store blocks}
        \If{$id, p \gets \Call{Learner.call\_set\_priority}$}
        \Comment{Get update priorities call}
            \State $\Call{set\_priority}{id, p}$
            \Comment{Update priorities}
    \EndIf
\EndFor
\EndProcedure
\end{algorithmic}
}
\end{algorithm}

\begin{algorithm}[ht]
\caption{RASPBERry Learner in Ape-X}
\label{alg:learner}
\small{
\begin{algorithmic}[1]
\Procedure{Learner}{$T$}\Comment{Train the network for $T$ updates using blocks from Replay}
    \State $\theta_0 \gets \Call{Initialize Network}{} $
    \For{$t = 1 \text{ to } T$}
        \State \textcolor{blue}{$(id, B_{\tau}) \gets \Call{ReplaySample}{}$}
        \Comment{Reconstruct original transitions in a training-ready format}
        \State $l_t \gets \Call{ComputeLoss}{\tau, \theta_t}$
        \Comment{Apply learning rule; e.g. double Q-learning or DDPG}
        \State $\theta_{t+1} \gets \Call{Update Parameters}{l_t; \theta_t}$
        \State $p \gets \Call{Compute Priorities}{\tau, \theta_{t+1}}$
        \Comment{Calculate priorities for experience, e.g., absolute TD error.}
        \State \textcolor{blue}{$B_{p} \gets \Call{Group priorities}{p}$}
        \Comment{Aggregate transition-level priorities into block-level priority}
        \State \textcolor{blue}{$\Call{Replay.Set\_Priority}{id, B_{p}}$}
        \Comment{Remote call to update priorities.}
        \State $\textproc{Periodically}(\Call{replay.Remove\_To\_Fit})$
        \Comment{Remove oldest blocks to make space}
\EndFor
\EndProcedure
\end{algorithmic}
}
\end{algorithm}


\section{Experiments setting}
\label{apx:experiments}
We provide details on the hardware and hyperparameters expanding on Section~\ref{sec:experiments} of the main paper.

\subsection{Hardware Details}
Experiments were conducted on two systems:
\begin{itemize}[leftmargin=*]
    \item \textbf{Local PC}: Intel 10900X CPU, NVIDIA RTX 3090 GPU, and 160\,GB DRAM.
    \item \textbf{HPC Cluster}: Multi-node setup with high-core-count CPUs and NVIDIA V100 GPUs.
    Unless otherwise stated, each HPC job utilized one V100 GPU and 20 CPU cores.\footnote{Exact cluster details are redacted for review; see \url{https://www.jade.ac.uk} for reference.}
\end{itemize}

\subsection{Hyperparameters for DDQN}
Table~\ref{tab:ddqn_hypers} summarizes the main hyperparameter values for our DDQN implementation. 
Most of hyperparameter match the default settings in Ray RLLib \citep{RayprojectRlexperiments_ray_2022}.
\begin{table}[h]
\centering
\begin{tabular}{l|l}
\toprule
\textbf{Hyperparameter} & \textbf{Value} \\
\midrule
Optimizer & Adam ($\epsilon = 1.5 \times 10^{-4}$) \\
Exploration Config & $\epsilon$ decays to $0.01$ over $2 \times 10^{5}$ steps \\
Learning Rate & $6.25 \times 10^{-5}$ \\
Training Start & $50,000$ steps \\
Batch Size & $32$ \\
Replay Capacity & $1 \times 10 ^{6}$ \\
Prioritized Replay $\alpha$ & $0.5$ \\
Prioritized Replay $\beta$ & $1.0$ \\
Target Network Update & every $8,000$ steps \\
$n$-step returns & $3$ \\
\bottomrule
\end{tabular}
\caption{Key hyperparameters for DDQN experiments.}
\label{tab:ddqn_hypers}
\end{table}
\subsection{Hyperparameters for Ape-X DDQN}
Table~\ref{tab:apex_hypers} lists hyperparameters used in Ape-X DDQN, adapted from \citet{DistributedPrioritizedExperience_horgan_2018} to fit our hardware constraints.

\begin{table}[h]
\centering
\begin{tabular}{l|l}
\toprule
\textbf{Hyperparameter} & \textbf{Value} \\
\midrule
Optimizer & Adam ($\epsilon = 1.5 \times 10^{-4}$) \\
Exploration Config & $\epsilon$ decays to $0.01$ over $200,000$ steps \\
$n$-step returns & $3$ \\
Learning Rate & $1.0 \times 10^{-4}$ \\
Training Start & $20,000$ steps \\
Batch Size & $512$ \\
Replay Capacity & $1 \times 10^{6}$ \\
Prioritized Replay $\alpha$ & 0.5 \\
Prioritized Replay $\beta$ & 1.0 \\
Target Network Update & every $5,000$ steps \\
Number of workers & $2$ \\
Number of environments per Worker & $4$ \\
$n$-step returns & $3$ \\
\bottomrule
\end{tabular}
\caption{Key hyperparameters for Ape-X DDQN experiments.}
\label{tab:apex_hypers}
\end{table}

\section{Is Compression Widely Applicable?}
\label{apx:compression_applicability}
Although our experiments focus mainly on image-based observations, RASPBERry’s compression stage can be adapted to varied data types.  
For instance, \texttt{zlib} \citep{MadlerZlib_adler_2024} achieves high compression on consecutive image frames, but it also can be used for different transition data.

Table~\ref{tab:ram_saved_diff} shows sample compression performance across diverse tasks.  
Even for non-image data, substantial reduction is possible.  
In practice, any compression library can be integrated, the best option may depend on the transitions' nature.

\begin{table}[h]
\centering
\begin{tabular}{llr}
\toprule
\textbf{Environment} & \textbf{Data Type} & \textbf{Compression Rate} \\
\midrule
Atari Breakout & Game Image & 88.90\% \\
Atari BattleZone & Game Image & 92.70\% \\
Atari Pong & Game Image & 95.70\% \\
Atari Breakout & Game ROM & 81.60\% \\
Atari Freeway & Game ROM & 73.28\% \\
Atari BeamRider & Game ROM & 81.04\% \\
Highway & Agent Metrics & 37.10\% \\
Pusher & Agent Metrics & 27.93\% \\
Reacher & Agent Metrics & 26.70\% \\
Airsim Car in Unity empty space* & Camera Image & 83.02\% \\
Airsim Drone in Unity FPS game* & Camera Image & 63.02\% \\
\bottomrule
\end{tabular}
\caption{Compression rates with \texttt{zlib} \citep{MadlerZlib_adler_2024} for different data types.
Block size $m = 32$.
\newline \textbf{*} Environments marked with `*' use 100$\times$100$\times$3 resolution.}
\label{tab:ram_saved_diff}
\end{table}

%%%%%%%%%%%%%%%%%%%%%%%%%%%%%%%%%%%%%%%%%%%%%%%%%%%%%%%%%%%%%%%%
%% NOTE: THIS MARKS THE END OF THE "MAIN TEXT"
%%%%%%%%%%%%%%%%%%%%%%%%%%%%%%%%%%%%%%%%%%%%%%%%%%%%%%%%%%%%%%%%

%%%%%%%%%%%%%%%%%%%%%%%%%%%%%%%%%%%%%%%%%%%%%%%%%%%%%%%%%%%%%%%%
%% Bibliography
%%%%%%%%%%%%%%%%%%%%%%%%%%%%%%%%%%%%%%%%%%%%%%%%%%%%%%%%%%%%%%%%
\bibliography{all_ref}
\bibliographystyle{rlj}

\medskip

%%%%%%%%%%%%%%%%%%%%%%%%%%%%%%%%%%%%%%%%%%%%%%%%%%%%%%%%%%%%%%%%
% AUTHOR: If your paper has no supplementary materials, you may 
%         comment out the line below, which creates the title for
%         the supplementary materials.
%%%%%%%%%%%%%%%%%%%%%%%%%%%%%%%%%%%%%%%%%%%%%%%%%%%%%%%%%%%%%%%%

\beginSupplementaryMaterials

\section{Supplementary Materials: HPC Experiments}
\label{sec:supplementary_hpc}
\subsection{Overview}
This section presents comprehensive experimental results obtained on a high-performance computing (HPC) cluster.\footnote{Exact cluster details are redacted for review; see \url{https://www.jade.ac.uk} for reference.}

\subsection{Metrics for HPC Experiments}
For completeness, we use those metrics in our evaluation:
\begin{itemize}[leftmargin=*]
    \item $R_{\text{improve}}$: The median percentage difference in the rolling-window average reward compared to the baseline.
    \item $T_{\text{save}}$: The median percentage of training time saved in reaching predefined reward thresholds.
    \item $R_{\text{max}}$: The highest episode reward recorded, expressed relative to the best reward of the baseline.
    The baseline is $100\%$
\end{itemize}

\subsection{HPC Experiment Summaries}
Table~\ref{tab:hpc_sum} compares \textbf{DDQN with RASPBERry} and \textbf{Ape-X~DDQN with RASPBERry} against their respective baselines.

\begin{table}[h]
\centering
\begin{tabular}{lrrrrr}
\toprule
 & $R_{\text{improve}}$ & $T_{\text{save}}$ & $R_{\text{max}}$ & Sampled & Trained \\
\midrule
DDQN + RASPBERry      &  60.4\% & 35.4\% & 102.9\%  & 160\% & 160\% \\
Ape-X DDQN + RASPBERry &  78.6\% & 56.5\% & 122.7\%  &  35\% & 1035\%  \\
\bottomrule
\end{tabular}
\caption{High-level comparison of RASPBERry vs.\ standard PER/DPER in HPC experiments.}
\label{tab:sm_hpc_sum}
\end{table}

Overall, these HPC trials confirm that RASPBERry’s block-based replay mechanism retains its benefits even at large scale system: 

\begin{itemize}[leftmargin=*]
    \item \(\mathbf{R_{\text{improve}}}\) and \(\mathbf{T_{\text{save}}}\) remain consistently high, indicating faster learning across most tasks.
    \item \(\mathbf{R_{\text{max}}}\) shows that final performance meets or surpasses the best baseline results.
    \item The `\emph{Sampled}' and `\emph{Trained}' columns reveal that RASPBERry processes substantially more frames in the same wall-clock time).
\end{itemize}

\subsection{Detailed HPC Experiment Result}
We additionally provide a more detailed comparison (Table~\ref{tab:hpc_detailed}) spanning 24 Atari environments.
The results further confirm that RASPBERry consistently improves training speed ($T_{\text{save}}$) and reward ($R_{\text{improve}}$), often reaching higher final scores ($R_{\text{max}}$) than PER/DPER.

It should be noted that we removed outliers in the Trimmed Mean, which may be due to the instability of training. 
However, large-scale experiments are expensive and we did not perform additional repeated tests, and it will not shake our conclusion.

\begin{table*}[htb]
\centering
% \footnotesize
\begin{tabular}{lcccccc}
\toprule
 & \multicolumn{3}{c}{\textbf{PER vs.\ RASPBERry}} & \multicolumn{3}{c}{\textbf{DPER vs.\ RASPBERry}} \\
\cmidrule(lr){2-4}\cmidrule(lr){5-7}
\textbf{Environment} & $R_{\text{improve}}$ & $T_{\text{save}}$ & $R_{\text{max}}$ & $R_{\text{improve}}$ & $T_{\text{save}}$ & $R_{\text{max}}$ \\
\midrule
Alien          & 65.7\%  & 69.3\%  & 130.1\% & 57.6\% & 67.5\% & 133.2\% \\
Amidar         & 20.2\%  & 50.0\%  & 101.5\% & 41.9\% & 57.2\% & 91.1\%  \\
Assault        & 60.9\%  & 85.2\%  & 111.8\% & 55.7\% & 85.2\% & 187.1\% \\
Asterix        & 1.3\%   & 0.0\%   & 102.3\% & 25.4\% & 71.4\% & 111.8\% \\
Asteroids      & 8.5\%   & 0.0\%   & 83.9\%  & -24.3\% & 0.0\% & 92.6\%  \\
Atlantis       & -2.1\%  & 0.0\%   & 100.0\% & -19.9\% & -175.0\% & 1447.4\% \\
BankHeist      & 720.0\% & 73.6\%  & 112.5\% & 0.0\%   & 0.0\% & 100.0\%  \\
BattleZone     & 72.7\%  & 95.8\%  & 134.4\% & 83.0\% & 83.0\% & 120.0\% \\
BeamRider      & 68.9\%  & 71.7\%  & 75.5\%  & 24.2\% & 73.1\% & 91.9\%  \\
Berzerk        & 55.8\%  & 0.0\%   & 87.4\%  & 20.9\% & 43.3\% & 135.9\% \\
Bowling        & 52.8\%  & 0.0\%   & 187.9\% & 98.3\% & 55.4\% & 107.0\% \\
Boxing         & 91.7\%  & 0.0\%   & 100.0\% & 29.5\% & 60.5\% & 100.0\% \\
Breakout       & 166.8\% & 82.7\%  & 104.4\% & 52.5\% & 88.1\% & 99.3\%  \\
Carnival       & -5.7\%  & 0.0\%   & 92.5\%  & 67.0\% & 46.8\% & 180.4\% \\
Centipede      & 43.1\%  & 52.6\%  & 109.6\% & 12.8\% & 0.0\% & 110.3\%  \\
ChopperCommand & -8.4\%  & 0.0\%   & 55.0\%  & 121.4\% & 71.1\% & 184.1\% \\
CrazyClimber   & 31.8\%  & 37.9\%  & 66.4\%  & 27.1\% & 56.1\% & 95.8\%  \\
Defender       & -7.7\%  & 0.0\%   & 93.5\%  & 14.8\% & 0.0\% & 189.7\% \\
FishingDerby   & 451.7\% & 81.9\%  & 303.8\% & 670.6\% & 87.3\% & 100.4\% \\
Freeway        & 36.7\%  & 88.4\%  & 103.0\% & 5634.7\% & 90.0\% & 106.2\% \\
Frostbite      & 65.6\%  & 50.0\%  & 84.9\%  & 32.8\% & 62.5\% & 119.8\% \\
Gopher         & -3.6\%  & -6.6\%  & 89.1\%  & 78.0\% & 79.1\% & 103.3\% \\
Qbert          & 45.7\%  & 27.7\%  & 101.2\% & 191.8\% & 88.0\% & 121.8\% \\
SpaceInvaders  & 8.2\%   & 8.2\%   & 92.1\%  & 42.9\% & 66.7\% & 108.3\% \\
\midrule
\textbf{Trimmed Mean} & 60.4\% & 35.4\% & 102.9\% & 78.6\% & 56.5\% & 122.7\% \\
\bottomrule
\end{tabular}
\caption{Comprehensive per-environment comparison for HPC experiments. Each metric is shown as a percentage relative to the PER or DPER baseline. The final row provides trimmed mean values to reduce the impact of outliers.}
\label{tab:hpc_detailed}
\end{table*}

\subsection{Experiment Results by Learning Curve}
To visualize performance trends over time, we plot learning curves (raw scores) vs.\ wall-clock time for each method. Figures~\ref{fig:jade_ddqn}~and~\ref{fig:jade_apex} show \emph{DDQN (RASPBERry vs.\ PER)} and \emph{Ape-X DDQN (RASPBERry vs.\ DPER)}, respectively.
All tasks used the same resource settings on the HPC cluster, and the curves are smoothed with a moving average for readability.

\begin{figure*}[ht]
\centering
\includegraphics[width=\textwidth]{images/experiments/all_hpc_per_time_10_paper.png}
\caption{Learning curves (raw score vs.\ wall-clock time) for \textbf{DDQN with RASPBERry vs.\ PER}, using a moving average smoothing.}
\label{fig:jade_ddqn}
\end{figure*}

\begin{figure*}[ht]
\centering
\includegraphics[width=\textwidth]{images/experiments/all_hpc_dper_time_10_paper.png}
\caption{Learning curves (raw score vs.\ wall-clock time) for \textbf{Ape-X DDQN with RASPBERry vs.\ DPER}, using a moving average smoothing    .}
\label{fig:jade_apex}
\end{figure*}

These findings highlight that the advantages of RASPBERry are not only preserved but can be amplified in a high-throughput, distributed training regime in HPC.
Its ability to compress and prioritize transitions at the block level appears to reduce cache misses and sum-tree updates, leading to higher overall efficiency.

\end{document}
